% ── Rozwiązanie: Exact Subset Sum -- instancja 2 ─────────────────
\subsection{Exact Subset Sum -- instancja 2}

Dane: $S = \{2, 3, 7, 8, 10\}$ (w~kolejności $x_1, \ldots, x_5$) oraz $t = 15$.
Algorytm \textsc{ExactSubsetSum} startuje z~$L_0 = (0)$ i~w~$i$-tej iteracji scala
$L_{i-1}$ z~$L_{i-1} + x_i$, po~czym usuwa elementy $> t$. Element wprowadzony
przez przesunięcie $L_{i-1} + x_i$ wyróżniono na~\textcolor{accentB}{pomarańczowo}.

\begin{dygresja}
	Przebieg iteracji ($L_{i-1} + x_i$ to składnik przesunięty):
	\begin{enumerate}[nosep]
		\item $x_1 = 2$:\quad
		      $(0) \cup (\textcolor{accentB}{2}) = (0, \textcolor{accentB}{2})$
		      \hfill $L_1 = (0, 2)$
		\item $x_2 = 3$:\quad
		      $(0, 2) \cup (\textcolor{accentB}{3}, \textcolor{accentB}{5})
		      = (0, 2, \textcolor{accentB}{3}, \textcolor{accentB}{5})$
		      \hfill $L_2 = (0, 2, 3, 5)$
		\item $x_3 = 7$:\quad scalenie z~$(7, 9, 10, 12)$, nic $> 15$
		      \hfill\\
		      $L_3 = (0, 2, 3, 5, \textcolor{accentB}{7}, \textcolor{accentB}{9},
		      \textcolor{accentB}{10}, \textcolor{accentB}{12})$
		\item $x_4 = 8$:\quad przesunięcie $L_3 + 8$ daje
		      $8, 10, 11, 13, 15, 17, 18, 20$; po~usunięciu $> 15$ dochodzą
		      $\textcolor{accentB}{8}, \textcolor{accentB}{11}, \textcolor{accentB}{13},
		      \textcolor{accentB}{15}$ ($10$ już było) \\
		      $L_4 = (0, 2, 3, 5, 7, \textcolor{accentB}{8}, 9, 10,
		      \textcolor{accentB}{11}, 12, \textcolor{accentB}{13}, \textcolor{accentB}{15})$
		\item $x_5 = 10$:\quad przesunięcie $L_4 + 10$ daje
		      $10, 12, 13, 15, 17, \ldots$; wszystkie $\leq 15$ ($10, 12, 13, 15$) już
		      są na liście, więc \emph{nic nowego}
		      \hfill $L_5 = L_4$
	\end{enumerate}
\end{dygresja}

\paragraph{Wynik.}
Największy element listy $L_5$ nieprzekraczający $t = 15$ to
\[
	\boxed{z = 15}
\]
czyli dokładnie $t$ -- istnieje podzbiór o~sumie równej progowi, np.\
$\{7, 8\}$ lub $\{2, 3, 10\}$. Wersja decyzyjna ma więc odpowiedź \emph{tak},
a~optymalna suma $\leq t$ wynosi $15$.
